\section{Finite evaluation of nested conditional means}\label{sec:compiler}

Suppose that a deterministic numerical formula calls for a conditional
mean at a center that depends on other conditional means.  Even a very
accurate unbiased estimate is awkward here: its error changes the next
center, and the next call must be analyzed at that random center.  We
will instead add a prescribed Gaussian to every computed center.  A
translation identity of the lower-level routine allows that Gaussian to
be absorbed into its random seeds.  The identity restores the law of a
call at the intended center, which is fixed after conditioning on the
parent formula's random input anchors.

We first prove that identity.  We then give the full finite recursive
construction, including its termination, its query cap on every
execution, and its distributional error.  The analytical input is a
family of finite affine formulas in exact conditional means.  Theorem~\ref{thm:numerical-interface} supplies these formulas and their
complete numerical construction.  We now implement them and thereby
prove Theorem~\ref{thm:primitive}.

\subsection{Why Gaussian center noise helps}

A seeded mean routine at scale $r>0$ is a measurable deterministic map
$\mathsf Q(x;z_1,\ldots,z_m)\in\R^d$, whose input is a center
$x\in\R^d$ and $m$ seed slots $z_j\in\R^d$.  It may query a fixed
oracle.  A vector $v=(v_1,\ldots,v_m)\in\R^m$ gives a
\emph{query-preserving translation} if, for every center, displacement
and fixed seed list,
\begin{equation}\label{eq:mean-shift-warmup}
 \mathsf Q(x+\Delta;z+v\otimes\Delta/r)=\mathsf Q(x;z),
\end{equation}
and the two computations ask exactly the same physical oracle queries.
Here $(v\otimes a)_j=v_ja$.  This is a deterministic identity; it does
not require independent or Gaussian seeds.

\begin{lemma}[Absorbing a Gaussian center displacement]
\label{lem:compiler-absorption}
Let $D_s>0$, and let $\mathsf Q$ satisfy
\eqref{eq:mean-shift-warmup}, with $\|v\|_2\le D_s$.  Put
\[
 n=\frac r{2D_s},\qquad
 P=\left(I_m-\frac{vv^T}{4D_s^2}\right)^{1/2}.
\]
For a deterministic $x\in\R^d$, let $Z,G_1,\ldots,G_m$ be independent
standard Gaussian vectors in $\R^d$.  Then the output of
$\mathsf Q(x+nZ;(P\otimes I_d)G)$ has exactly the law of
$\mathsf Q(x;G)$.
\end{lemma}
\begin{proof}
The eigenvalues of $P$ lie in $[\sqrt3/2,1]$.  The translation identity
rewrites the first output as
\[
 \mathsf Q\left(x;PG-\frac{v\otimes Z}{2D_s}\right).
\]
We suppress the tensor product with $I_d$ on slot-space matrices.
The seed array in this expression is centered Gaussian and has slot
covariance
\[
 PP^T+\frac{vv^T}{4D_s^2}=I_m.
\]
Its slots are therefore independent standard Gaussian vectors.  They
may remain correlated with $Z$; that is harmless because the rewritten
call is at the deterministic center $x$.
\end{proof}

The covariance reduction $P$ is essential.  With unmodified fresh
seeds, translating the noisy center back would add variance.  With
$P$, the center noise restores precisely the variance removed from the
seeds.  The construction below maintains the deterministic translation
identity at every level, so that this argument applies recursively.

\subsection{The numerical input and the implementation task}

We retain the primitive notation of Section~\ref{sec:preliminaries}:
$g=\nabla F$, $\operatorname{Lip}(g)\le\lambda=\delta=d^{-e}$,
$0<r\le1$, and $\nu_{x,r}$ is the law in
\eqref{eq:primitive-law}.  Write $u_{x,r}=\E_{\nu_{x,r}}g(x+rZ)$,
$L=\lambda r$, $l=\lambda r^2$, and $D(x)=\sqrt d+|g(x)|$.
We use the uniform moment bounds \eqref{eq:primitive-moments} and
conditional coupling rules of Lemma~\ref{lem:coupling-smoothing}.
Only evaluations of $g$ will be used by the implementation.

Fix $K_0,t$ as in Theorem~\ref{thm:primitive}, and retain the
numerical parameters
\[
 R=\delta^t,\qquad b_*=(1-R^2)^{1/2},\qquad
 \psi=\delta^{4t},\qquad J=K_0+2.
\]
Theorem~\ref{thm:numerical-interface} supplies two finite acyclic
circuits, $\mathcal C_M$ and $\mathcal C_S$, whose nodes are exact
conditional means.  We use its notation $A$ for the positive declared
bound on a center's absolute coefficient sum, $A_f$ for the output
bound, $D_0$ for the mean-dependency depth, and $c_1$ for the absolute
constant in the fan-out exponent.  A center has the form
$x+r(\ell_i y+m_iG)+\sum_{j<i}a_{ij}u_{x_j,r_j}$, with
$r_i=r\tau_i$.  All these scalar data are deterministic functions of
the prescribed numerical parameters, independent of centers, anchors,
and oracle responses.  We will use the circuit theorem's explicit
coefficient and scale bounds when constructing shifts and propagating
sizes; the circuits themselves are comparison formulas, not additional
runtime oracles.

For the mean circuit the parameter supplied is the parent's
$\sigma$.  Its underlying centering flow uses $s=\sigma/2$ once,
so its direct output is $(\sigma/2)G$.  Under the specified ideal
anchor law, its error is
\[
 W_2\bigl(\mathcal L(\mathcal C_M),N(u_{x,r},\sigma^2I_d/4)\bigr)
 \le\polylog(d)\sigma\delta^JD(x).
\]
For the sampling circuit the direct
output is $y/\rho_e$, where
$\rho_e=(1+\eta^2/4)^{-1/2}$, and the ideal target is
$\nu_{x,r}*N(0,\eta^2I_d/4)$.  These numerical statements require
the anchor laws and size assumptions specified in
Theorem~\ref{thm:numerical-interface}; we shall verify those assumptions
at their uses.

The remaining task is to replace every exact mean by finitely many
queries while preserving enough independent Gaussian noise for the
absorption identity.  The construction below first establishes a
query cap on every execution.  Its subsequent accuracy proof uses
pathwise stability at inaccurate centers, then absorption at ideal
noisy centers, and only then a standalone law for a child call.

\subsection{Two routines and two different kinds of depth}

We construct a mean routine $\mathsf M(x,r,\sigma;b,p)$ targeting
$u_{x,r}+\sigma G$ and a sampling routine
$\mathsf S(x,r,\eta;b,p)$ targeting $Z+\eta G$, for independent
$Z\sim\nu_{x,r}$ and standard Gaussian $G$.
The top call is $\mathsf S(x,r,1;1,K_0)$.  Throughout, a reachable
call means a descendant of this call, including the call itself.
The first label $b$ records the size conditions
\begin{equation}\label{eq:routine-sizes}
 \max\{l,L/\sigma\}\le\polylog(d)\delta^b\quad(\mathsf M),
 \qquad l/\eta\le\polylog(d)\delta^b\quad(\mathsf S),
\end{equation}
where $R\le\eta\le1$.  The second label $p$ records the requested
error power.  The intended errors are
$\polylog(d)\sigma\delta^pD(x)$ and
$\polylog(d)\delta^pD(x)$, respectively.
A call is terminal precisely when $p\le b$; it is nonterminal
when $p>b$.
Every sampling routine reserves a final noise $\eta G/2$, unused
in its own queries.  Before adding that reserve, its target law is
$\nu_{x,r}*N(0,3\eta^2I_d/4)$.

For a nonterminal sampling call, use $\mathcal C_S$ with a standard
Gaussian position anchor, supply additional final noise
$n=\eta/\sqrt2$, and then add the reserve.  For a nonterminal mean
call, pass the parent parameter $\sigma$ to $\mathcal C_M$ and supply
final noise $n=\sqrt3\sigma/2$.  Obtain its position anchor from an
$\mathsf S$ call at the same center and scale with
$\eta=R/b_*$, multiplying that output by $b_*$.  An ideal such
anchor has exactly the required law $b_*Z+RN$.
Its momentum anchor is an independent standard Gaussian.  For all
sufficiently large $d$, $R/b_*\in[R,1]$, so the anchor call has an
admissible noise parameter.

An expression dependency and a routine call play different roles.
For example, suppose a parent circuit contains
$E_f=a\,u_{x_2,r_2}$ and $x_2=x+c\,u_{x_1,r_1}$.
Evaluating $E_f$ invokes two new mean routines: one to help supply
$x_2$, and one at that supplied center.  Both are children of the
\emph{same parent routine}.  Their own internal circuits begin the
next routine level.  The two nested affine expressions do not create
two successive losses in the parent's accuracy label.  Figure
\ref{fig:two-depths} makes this distinction explicit.

\begin{figure}[tbp]
\centering
\begin{tikzpicture}[>=Latex,font=\small,
 box/.style={draw,rounded corners,align=center,minimum width=36mm,minimum height=12mm}]
 \node[box] (p) at (0,1.7) {parent routine};
 \node[box] (a) at (-2.8,0) {mean routine\\at $x_1$};
 \node[box] (b) at (2.8,0) {mean routine\\at supplied $x_2$};
 \draw[->] (p)--(a);\draw[->] (p)--(b);
 \draw[->,dashed] (a.south) -- (-2.8,-1) --
 node[below] {supplies center} (2.8,-1) -- (b.south);
 \node[align=center] at (0,-1.8) {solid: routine parentage\qquad dashed: expression dependence};
\end{tikzpicture}
\caption{The two calls have the same routine parent even though
one helps form the other's center.  When a mean occurs twice, each
occurrence uses a separate copy of its entire center-supplying subtree.}
\label{fig:two-depths}
\end{figure}



To evaluate an expression, we expand its mean dependencies as an
uncached tree, duplicating every mean occurrence together with its
entire center-supplying subtree.  All mean calls at any expression tier
are direct children of the routine whose circuit is being evaluated.

The children of a parent labeled $(b,p)$ receive labels as follows:
\begin{equation}
 \begin{array}{c|c}
 \text{child's role}&\text{child labels}\\ \hline
 \text{term of }E_f\text{ in an }\mathsf M\text{ parent}&(b+t,p)\\
 \text{term of }E_f\text{ in an }\mathsf S\text{ parent}&(b,p)\\
 \text{any other term, or the }\mathsf M\text{ position anchor}
       &(b-10t,p-1/4).
 \end{array}
 \label{eq:compiler-labels}
\end{equation}
In the last row the labels are always those of the parent routine,
not successively decreased while descending its expression tree.

\begin{lemma}[Bounded routine height]
\label{lem:compiler-height}
Starting at $(b,p)=(1,K_0)$, the rules
\eqref{eq:compiler-labels} give a finite routine tree of height at most
a constant $H=H(K_0,t)$.  Every reachable label satisfies $b>1/2$.
\end{lemma}
\begin{proof}
Before $4K_0+1$ decreasing transitions, the total loss in $b$ is less
than $10t(4K_0+1)<1/4$, so $b>3/4$.  At that many transitions,
$p\le K_0-(4K_0+1)/4<0$, forcing termination.  Therefore no branch
can have more such transitions.  Positive transitions increase $b$
by $t$, while a nonterminal satisfies $b<p\le K_0$; their total
number is consequently bounded in terms of $K_0,t$ and the already
bounded total decrease in $b$.  A neutral transition is an
$\mathsf S$-to-$\mathsf M$ transition.  Another neutral transition
requires first reaching another $\mathsf S$, which can occur only
through a decreasing transition.  This bounds the neutral transitions
as well and proves both claims.
\end{proof}

A routine is a deterministic function of its center, scalar
parameters, oracle, and a finite ordered list of $d$-dimensional seed
slots.  Standalone use fills these slots by independent standard
Gaussians.  In addition, each routine will have a deterministic vector
$v$ of scalar coefficients, one per slot, with
\begin{equation}
 \|v\|_2\le D_s,\qquad D_s:=\delta^{-t}.
 \label{eq:shift-budget}
\end{equation}
For \emph{every} fixed seed list, changing $x$ to $x+\Delta$ and
adding $v\otimes\Delta/r$ to the seeds will preserve every physical
query point.  The output of $\mathsf M$ will be unchanged; the
normalized output of $\mathsf S$ will decrease by $\Delta/r$.
These are pathwise properties, not statements restricted to Gaussian
seeds.

Within each call write $u=u_{x,r}$ and $\nu=\nu_{x,r}$.
At a terminal call define
\begin{align*}
 \mathsf M&=g(x+rZ)+\sigma G,\\
 \mathsf S&=Z-rg(x+rZ)+(\sqrt3\eta/2)G_1+(\eta/2)G_2,
\end{align*}
with separate seed slots.  The last term in $\mathsf S$ is its
reserved noise.  Taking $v=-1$ on the $Z$ slot and zero elsewhere
proves the shift property and $\|v\|_2=1$.



\subsection{Evaluating an expression at a noisy center}

Fix a parent circuit and its anchors, and write $M_i=u_{x_i,r_i}$
for its exact mean nodes.  For each expression
$E=b_0+\sum_i a_iM_i$, the operation $\operatorname{Eval}(E,n)$
will approximate its exact value plus an independent
$N(0,n^2I_d)$ noise.  Assign disjoint seed blocks to every term
occurrence, its center-supplying computation, and each residual noise;
keep all these blocks separate from the anchors.  Thus the term
blocks are independent conditional on the anchors in a standalone
call.

Here is the recursive evaluation rule, assuming the child routines
and their shift vectors have already been constructed.  Take $A'=A$
for a center expression and $A'=A_f$ at the output, and put
$\sigma_i=n/(2A')$.  For the term of scale $r_i=r\tau_i$, first
evaluate its center expression with noise
\begin{equation}
 n_i=\frac{r_i}{2D_s}.
 \label{eq:compiler-center-noise}
\end{equation}
At this supplied center run the child $\mathsf M$ with scale $r_i$,
noise $\sigma_i$, and labels from \eqref{eq:compiler-labels}.
If its shift vector is $v_i$, apply to its fresh seed block $G_i$
the deterministic slot-space transformation
\begin{equation}
 P_i=\left(I-\frac{v_iv_i^T}{4D_s^2}\right)^{1/2},
 \qquad\text{and use seeds }(P_i\otimes I_d)G_i.
 \label{eq:compiler-seed-transform}
\end{equation}
Denote its output by $Q_i$.  Return
\begin{equation}
 b_0+\sum_i a_iQ_i+
 \left(n^2-\sum_i a_i^2\sigma_i^2\right)^{1/2}G_E,
 \label{eq:compiler-eval}
\end{equation}
where $G_E$ is a separate slot.  The square root is defined, since
\[
 \sum_i a_i^2\sigma_i^2
 \le\left(\sum_i|a_i|\right)^2\frac{n^2}{4(A')^2}
 \le n^2/4.
\]
If there are no terms, return $b_0+nG_E$.  The mean-node ordering
makes the expression recursion finite.  At the parent's output use
the value of $n$ specified above; add the reserved noise for
$\mathsf S$.  A routine supplied with correlated seeds performs
these same deterministic operations, without any change of rule.

At an ideal supplied center $x_i+n_iZ$, Lemma~\ref{lem:compiler-absorption}
now identifies this child with its standalone law at $x_i$.  We still
owe the construction of the shift vectors and a proof that the noisy
supplied centers are accurate.  We prove these in that order.

\subsection{Query-preserving shifts and the deterministic cap}

We next show that the shift vectors required in
Lemma~\ref{lem:compiler-absorption} exist with the universal bound
\eqref{eq:shift-budget}.  This also verifies that the definition is
noncircular and has a deterministic query cap.

For a nonterminal $\mathsf S$, shift its Gaussian position anchor by
$-\rho_e\Delta/r$.  For a nonterminal $\mathsf M$, we require its
position anchor $y_*$ to shift by $-\Delta/(b_*r)$ and leave its
momentum anchor unchanged.  To obtain this, apply the shift of its anchor
$\mathsf S$ call and additionally shift that call's reserved seed
by the coefficient $-2R/b_*$ times $\Delta/r$.  The ordinary shift
changes its output by $-\Delta/r$.  Since the reserved standard
deviation is $R/(2b_*)$, the additional change is
$-R^2\Delta/(b_*^2r)$.  Multiplication by $b_*$ makes the total
change
\[
 -\left(b_*+\frac{R^2}{b_*}\right)\frac{\Delta}{r}
 =-\frac{\Delta}{b_*r}.
\]
The extra slot is used in no query, so query invariance is preserved.

For a term of scale $r\tau_i$, the direct part of its supplied
center then changes by $k_i\Delta$, where
\[
 k_i=1-\ell_i\rho_e\quad(\mathsf S),
 \qquad
 k_i=1-\ell_i/b_*\quad(\mathsf M).
\]
These coefficients satisfy $|k_i|/\tau_i\le C$.  For $\mathsf S$,
$\rho_i\le\rho_e$ gives
$1-\rho_i\rho_e\le2(1-\rho_i)\le2\tau_i$.
For $\mathsf M$, its two possible coefficients obey
\[
 0\le1-\rho_i\cos z/b_*
 \le C(1-\rho_i)+Cz^2,
\]
with $z=0$ when there is no cosine.  Use
$1-\rho_i\le\tau_i$, $\tau_i\ge cR$, and
$z^2/R\le C\delta^{7t}\to0$.

Shift that child's untransformed seed block by
\begin{equation}
 \frac{k_i}{\tau_i}P_i^{-1}v_i\otimes\frac{\Delta}{r}.
 \label{eq:compiler-child-shift}
\end{equation}
If \eqref{eq:shift-budget} holds for the child, the smallest
eigenvalue of $P_i$ is at least $\sqrt3/2$, so
$\|P_i^{-1}\|_{\op}\le2$.  Formula
\eqref{eq:compiler-child-shift} induces exactly the child's shift
for center displacement $k_i\Delta$.  Leave all residual noise
slots fixed.  Induction from leaf expressions upward shows that
every mean output is unchanged, and therefore that the actual
supplied-center change is precisely its direct-term change.
All queries are preserved.  The parent's final direct term is
unchanged for $\mathsf M$ and changes by $-\Delta/r$ for
$\mathsf S$, giving the required output behavior.

The center-supplying subtree and the invoked child have disjoint
slots.  In particular, the factor $k_i/\tau_i$ in
\eqref{eq:compiler-child-shift} multiplies only the shift of that
child, not the shifts of routines used to supply its center.  Those
are other direct children of the same parent and are charged once
in its total shift vector.

\begin{lemma}[Well-defined construction and query cap]
\label{lem:compiler-count}
One can choose $w>0$, depending only on $K_0,t$ and the fixed Picard
counts, and then choose the interpolation orders, so that the entire
construction is well defined for all sufficiently large $d$.
Every routine reached from a top call satisfies
\eqref{eq:shift-budget}.  Its total number of seed slots and oracle
queries is at most $\polylog(d)\delta^{-t}$ on every execution.
\end{lemma}
\begin{proof}
Let $D_0$ be the dependency-depth bound in
Theorem~\ref{thm:numerical-interface}, and let $F_0$ be its direct
fan-out bound.  Expanding a circuit into an expression tree costs
at most $(D_0+1)(1+F_0)^{D_0}$ nodes.  Its size is thus
$\polylog(d)\delta^{-C_2w}$, where $C_2$ depends on the Picard
counts but not on $w$ or the interpolation orders.  This bounds its
direct child invocations and its own anchor and residual slots.
Expanding through the height $H$ from
Lemma~\ref{lem:compiler-height} gives
$\polylog(d)\delta^{-C_3w}$ slots and terminal queries.
Interpolation orders change only fixed multiplicative constants in
these bounds, not $C_2$ or $C_3$.

The terminal shift has norm $1$.  Assuming the shift bound for the
children, the bounded ratios and inverse norms above show that a
parent's shift norm is at most a fixed constant times one plus the
sum of its direct children's shift norms, including the anchor
call.  The extra reserved-slot shift has bounded coefficient.
The triangle inequality is sufficient even if ordinary and extra
anchor shifts affect the same slot.  Iterating through the fixed
routine height gives the same bound
$\polylog(d)\delta^{-C_3w}$, after increasing $C_3$ if necessary.

Choose $w>0$ with
\[
 C_3w<t/2,\qquad w<t,\qquad c_1w<t,
\]
and only then choose the fixed interpolation orders needed in
Theorem~\ref{thm:numerical-interface}.  For sufficiently large $d$, the last
shift bound is smaller than $D_s=\delta^{-t}$.  This is a valid
bottom-up induction: it proves each child's bound before using its
inverse transformation in the parent.  The count bound is also at
most $\polylog(d)\delta^{-t}$.

To make the parameter order explicit, first fix the Picard counts,
hence $D_0,H,C_3$, then $w$, then the interpolation orders.  The
universal scalar $D_s$ fixes every center-noise budget.  The circuit
shapes, labels, requested noises, and slot counts are determined by
these scalar parameters before any shifts are computed.  Compute
child shifts before their matrices and the parent shift.  No center
value, oracle response, numerical-error estimate, or assumption
\eqref{eq:routine-sizes} enters this finite construction or the
count.  Only terminal routines query $g$, once each; hence the bound
is a cap on every execution, not an expected query count.
\end{proof}

\subsection{Terminal errors and propagation of scales}

The finite construction and its cap are now established independently
of accuracy.  We next verify its terminal errors and the size conditions
needed by the numerical estimates.

\begin{lemma}[Terminal accuracy]
\label{lem:compiler-terminal}
In standalone use, a terminal call satisfying \eqref{eq:routine-sizes}
obeys
\begin{align*}
 W_2\bigl(\mathcal L(\mathsf M),N(u,\sigma^2I_d)\bigr)
 &\le\polylog(d)\sigma\delta^pD(x),\\
 W_2\bigl(\mathcal L(\mathsf S_{\rm pre}),
          \nu_{x,r}*N(0,3\eta^2I_d/4)\bigr)
 &\le\polylog(d)\delta^pD(x),
\end{align*}
where $\mathsf S_{\rm pre}$ is the output before adding its reserved
noise $\eta G_2/2$.  The sampling bound also holds after adding that
same independent reserve to both laws.  The fixed-seed center
Lipschitz constant is at most $\lambda$ for $\mathsf M$ and
$r\lambda$ for $\mathsf S$.
\end{lemma}
\begin{proof}
For the standard Gaussian seed $Z$, the primitive moment bounds give
\[
 \|g(x+rZ)-u\|_{L^2}
 \le CL\bigl(\sqrt d+r|g(x)|\bigr)\le CLD(x).
\]
Since $L/\sigma\le\polylog(d)\delta^b$ and $p\le b$, this is
within the mean error budget; couple the added Gaussian identically.

Let $z_0$ be the mode of $\nu_{x,r}$, so
$z_0=-rg(x+rz_0)$.  Then
\[
 \|Z-rg(x+rZ)-(Z+z_0)\|_{L^2}
 \le l\|Z-z_0\|_{L^2}\le ClD(x).
\]
Also $W_2(N(z_0,I_d),\nu_{x,r})\le Cl\sqrt d$.
Apply the synchronous comparison in
Lemma~\ref{lem:convex-preliminaries} with Gaussian potential
$H_0(z)=|z-z_0|^2/2$ (contraction constant $1$) and tilted potential
$H(z)=|z|^2/2+F(x+rz)$.  Their gradient difference is
$r[g(x+rz)-g(x+rz_0)]$, whose $L^2(\nu_{x,r})$ norm is at most
$l\|Z-z_0\|_{L^2(\nu_{x,r})}\le Cl\sqrt d$.
This proves the displayed Wasserstein bound under exactly the
$C^2$ and bounded-Hessian hypotheses of that lemma.
Adding the same independent comparison noises proves the sampling
error $ClD(x)\le\polylog(d)\eta\delta^bD(x)$, which is at most
the requested bound because $\eta\le1$ and $p\le b$.  Omitting
$G_2$ proves the pre-reserved assertion.  The Lipschitz bounds follow
directly from the two displayed terminal formulas.
\end{proof}

\begin{lemma}[Size conditions descend through the routine tree]
\label{lem:compiler-sizes}
Starting with a top call satisfying
\eqref{eq:routine-sizes}, every descendant satisfies that condition
with its assigned labels, up to a uniform polylogarithmic factor.
At each nonterminal circuit,
\begin{equation}
 \lambda A\le\polylog(d)\delta^{1/3}.
 \label{eq:small-coefficients}
\end{equation}
For nonfinal expression terms within an $\mathsf M$ parent,
$\sigma_i\le C\sigma$; within an $\mathsf S$ parent,
$r\sigma_i\le C$.
\end{lemma}
\begin{proof}
For a final term of an $\mathsf M$ parent, the output uses only means of scale $rR$, so $r_i=rR$, while $\sigma_i$ is a fixed positive multiple
of $\sigma$.  Thus $L_i/\sigma_i$ gains a factor $R$ and $l_i$
gains $R^2$, allowing labels $(b+t,p)$.
For a final term of an $\mathsf S$ parent,
$\sigma_i$ is a fixed positive multiple of $\eta/r$.
Therefore $L_i/\sigma_i\le Cl/\eta$ and $l_i\le l$, allowing
the unchanged labels.

For a nonfinal term, let $r\tau_{\rm prev}$ be the scale of the
node whose center is currently being supplied.  Its expression is
evaluated with budget $r\tau_{\rm prev}/(2D_s)$, irrespective of
the tier at which it occurs.  Consequently
\[
 \sigma_i=\frac{r\tau_{\rm prev}}{4D_sA},
 \qquad
 \frac{L_i}{\sigma_i}
 \le \frac{CD_s\lambda A}{R}.
\]
The coefficient estimates give
$\lambda A\le\polylog(d)\delta^{b-4t-c_1w}$ for an
$\mathsf M$ parent.  For an $\mathsf S$ parent,
$\lambda A=C_1l\le\polylog(d)\delta^b$, since $\eta\le1$.
Using $D_s/R=\delta^{-2t}$ and $c_1w<t$ costs less than the
allowed $10t$ decrease in $b$.  Also $l_i\le l$.
For the position-anchor $\mathsf S$ call of $\mathsf M$,
$l/\eta\le l/R$, which has the same allowance.
This proves the size conditions by the bounded routine height.
All accumulated fixed and polylogarithmic factors are still one
polylogarithmic factor.

As $b>1/2$ and $4t+c_1w<5t<1/6$, the coefficient bounds imply
\eqref{eq:small-coefficients}.  Finally, the positive declared bound
for $A$ includes a fixed multiple of $r/\sigma$ in an
$\mathsf M$ parent, so
$\sigma_i\le r/(4D_sA)\le C\sigma$.
For $\mathsf S$, $A=C_1r^2$ gives
$r\sigma_i\le1/(4C_1D_s)\le C$.
\end{proof}

\subsection{Stability for arbitrary fixed seeds}

\begin{lemma}[Fixed-seed stability]
\label{lem:compiler-lipschitz}
For every seed array, including correlated or transformed arrays,
the compiled $\mathsf M$ output is a function of $x$ with Lipschitz
constant at most $\polylog(d)\lambda$; the normalized
$\mathsf S$ output has constant at most
$\polylog(d)r\lambda$.  If the anchors of a nonterminal
$\mathsf M$ are treated as separate inputs, its Lipschitz constant
in the position anchor is at most $\polylog(d)L$.
\end{lemma}
\begin{proof}
Induct on routine height, with the terminal case supplied by
Lemma~\ref{lem:compiler-terminal}.  First hold the anchors fixed.
A direct center term has Lipschitz constant $1$ in $x$ and at most
$Cr$ in a position anchor.  The center-supplying expression combines
child outputs with total absolute coefficient at most $A$.
Using the child $\mathsf M$ bound, the contribution of earlier
supplied-center sensitivities is at most
$\polylog(d)\lambda A$ times their maximum.  By
\eqref{eq:small-coefficients}, this is small.  Induction through
the bounded expression depth therefore bounds center sensitivities
by a constant times their direct sensitivities, with an allowed
polylogarithmic factor.

At the final expression the coefficient sum is $A_f$.  Its direct
term has no center dependence when anchors are held fixed; for
$\mathsf M$ it also has no position-anchor dependence.  Hence its
center sensitivity is at most
$\polylog(d)A_f\lambda$, and its $\mathsf M$ position-anchor
sensitivity at most $\polylog(d)A_f\lambda r$.
These are the asserted bounds with anchors treated as separate inputs.
The $\mathsf S$ Gaussian anchor is indeed center independent.
The $\mathsf M$ position anchor is supplied by $\mathsf S$ at
the same center, whose center Lipschitz constant by induction is
$\polylog(d)r\lambda$.  Its additional contribution is at most
$\polylog(d)Lr\lambda=\polylog(d)L^2$, and
$L^2\le\lambda$ for $r\le1$ and $\lambda\le1$.
All seed transformations depend only on scalar parameters, not on
$x$.  The proof uses only fixed-seed identities and Lipschitz
bounds, and therefore requires no seed independence.
\end{proof}

\subsection{Standalone accuracy and the primitive theorem}

We prove the law guarantees by induction on routine height.  This is
the only stage at which a routine's seeds are assumed independent
standard Gaussians.  Such an assumption will be used for a child only
after the absorption identity has restored its full standalone seed
law.

\begin{proposition}[Accuracy of the compiled routines]
\label{prop:compiler-accuracy}
For every call reached from the top labels, the standalone routines
satisfy
\begin{align*}
 W_2(\mathcal L(\mathsf M),N(u,\sigma^2I_d))
 &\le\polylog(d)\sigma\delta^pD(x),\\
 W_2\bigl(\mathcal L(\mathsf S_{\rm pre}),
       \nu_{x,r}*N(0,3\eta^2I_d/4)\bigr)
 &\le\polylog(d)\delta^pD(x),
\end{align*}
where $\mathsf S_{\rm pre}$ denotes the output before its reserved
noise.  The latter bound also holds after adding the reserved noise
to both laws, and
\[
 \|\mathcal L(\mathsf S)-\nu_{x,r}*N(0,\eta^2I_d)\|_{\TV}
 \le\polylog(d)\frac{\delta^pD(x)}{\eta}.
\]
\end{proposition}

\begin{proof}
The terminal assertion is Lemma~\ref{lem:compiler-terminal}.
All random variables below have finite second moments: the actual
algorithms are finite compositions of affine maps and globally
Lipschitz oracle evaluations, while ideal tilted variables have the
primitive moment bounds.  Thus all $W_2$ comparisons are defined.

For a nonterminal $\mathsf M$, first couple its approximate
$\mathsf S$ position anchor to an ideal anchor, independently of
all remaining seed blocks.  The child labels are
$(b-10t,p-1/4)$.  By induction and
Lemma~\ref{lem:compiler-lipschitz}, replacing this anchor costs at
most
\[
 \polylog(d)L\delta^{p-1/4}D(x).
\]
After division by the mean error unit $\sigma\delta^pD(x)$,
the cost is at most
$\polylog(d)\delta^{b-1/4}$, hence is bounded for large $d$.
We may now assume that the position anchor has law $b_*Z+RN$,
where $Z\sim\nu_{x,r}$ and $N$ is standard Gaussian, and that the
momentum anchor is an independent standard Gaussian.  These are exactly
the input laws in Theorem~\ref{thm:numerical-interface}.

Let $x_i$ be the centers of the deterministic exact-mean circuit
under these anchors, and put $D_i=\sqrt d+|g(x_i)|$.
Their RMS values over anchors obey
\begin{equation}
 \max_i\|D_i\|_{L^2}\le\polylog(d)D(x).
 \label{eq:compiler-center-moments}
\end{equation}
Indeed, both position and momentum anchors have RMS at most $CD(x)$.
Use $|g(x_i)|\le|g(x)|+\lambda|x_i-x|$ and the circuit's
direct terms.  An exact primitive mean at a preceding center has
magnitude at most $CD_j$.  Their contribution to this maximum-of-RMS
recurrence is at most $C\lambda A$ times its previous maximum.
This is small by \eqref{eq:small-coefficients}; the direct terms cost
at most $CD(x)$.  Induction through the node ordering proves
\eqref{eq:compiler-center-moments} without a maximum inside the
expectation.

Condition on the anchors.  For an expression $E$, write
$\epsilon_E$ for the $W_2$ error of
$\operatorname{Eval}(E,n)$ relative to its deterministic exact
value plus $N(0,n^2I_d)$.  Couple the supplied center of its $i$th
term to $x_i+n_iZ$ with its conditional $W_2$ error
$\epsilon_{E_i}$.  Use the same transformed child seeds on these
two centers, drawn independently of this coupling.  The fixed-seed
Lipschitz bound costs at most
$\polylog(d)\lambda\epsilon_{E_i}$.
At the ideal noisy center, Lemma~\ref{lem:compiler-absorption}
identifies the child's law with its standalone law at $x_i$.
The induction hypothesis therefore supplies the additional error
$\polylog(d)\sigma_i\delta^{p_i}D_i$.

Different term occurrences have disjoint blocks, so their outputs
are independent conditional on the anchors.  Product couplings
compare them to independent noised exact means.  Together with the
independent residual in \eqref{eq:compiler-eval}, the ideal sum
has exactly covariance $n^2I_d$.  We conclude that
\begin{equation}
 \epsilon_E\le\sum_i|a_i|
 \left[\polylog(d)\lambda\epsilon_{E_i}
       +\polylog(d)\sigma_i\delta^{p_i}D_i\right].
 \label{eq:compiler-error-recursion}
\end{equation}
Notice the order of this argument: inaccurate supplied centers use
only pathwise stability, and standalone accuracy is invoked only
after Gaussian absorption.  No assertion is made that the actual
algorithm's intermediate states have the ideal input laws.

Take RMS in the anchors and expand
\eqref{eq:compiler-error-recursion} through the finite expression
tree.  At its root, direct-child labels have $p_i=p$ and
\[
 \sum_i|a_i|\sigma_i
 \le A_f\frac{n}{2A_f}=n/2.
\]
Thus their intrinsic errors fit the requested budget.  At tier
$j\ge2$, absolute coefficient sums and Lipschitz factors contribute
at most
\[
 A_f\bigl(\polylog(d)\lambda A\bigr)^{j-1}.
\]
All such children have the single accuracy-label loss $p_i=p-1/4$,
not one loss per expression tier.  Their noise units are at most
$C\sigma$ in an $\mathsf M$ parent and $C/r$ in an
$\mathsf S$ parent by Lemma~\ref{lem:compiler-sizes}; these units
combine with $A_f=C_1$ or $C_1r$, respectively.
Already the first factor
$\polylog(d)\lambda A\le\polylog(d)\delta^{1/3}$
pays for the loss $\delta^{-1/4}$, leaving a positive power
$\delta^{1/12}$.  The bounded expression depth allows all further
tiers and polylogarithmic factors to be summed.  Leaf affine
expressions have zero error.  This proves the required final
evaluation error.

Conditionally on the anchors, the ideal final evaluation is its
deterministic exact-mean circuit value plus a Gaussian of the
specified fixed covariance.  Its noise may therefore be taken
independent of the anchors.  By Lemma~\ref{lem:compiler-sizes},
Theorem~\ref{thm:numerical-interface} applies.  It compares the ideal mean
circuit to $N(u,\sigma^2I_d/4)$ and the ideal sampling circuit to
$\nu_{x,r}*N(0,\eta^2I_d/4)$, with the respective errors
$\polylog(d)\sigma\delta^JD(x)$ and
$\polylog(d)\delta^JD(x)$.  Since $J=K_0+2$ and $p\le K_0$,
both errors fit the requested bounds.
Couple the independent evaluation noise identically in each comparison.
For $\mathsf M$, its variance $3\sigma^2/4$ raises the Gaussian
target variance from $\sigma^2/4$ to $\sigma^2$.
For $\mathsf S$, its variance $\eta^2/2$ raises the Gaussian
convolution variance from $\eta^2/4$ to $3\eta^2/4$.
This closes the standalone induction.  Adding the common reserved
noise preserves the $W_2$ bound; Gaussian smoothing by its standard
deviation $\eta/2$ gives the stated total variation bound by
\eqref{eq:gaussian-smoothing}.
\end{proof}

\begin{proof}[Proof of Theorem~\ref{thm:primitive}]
Use the top $\mathsf S$ call with $(b,p,\eta)=(1,K_0,1)$.
Because $l=\delta r^2\le\delta$, its size condition holds with
constant $1$.  Proposition~\ref{prop:compiler-accuracy} gives both
$W_2$ and total variation error at most
$\polylog(d)\delta^{K_0}D(x)$ from
$\nu_{x,r}*N(0,I_d)$, uniformly in $x$ and $0<r\le1$.
Lemma~\ref{lem:compiler-count} gives at most
$\polylog(d)\delta^{-t}$ evaluations of $g$ on every execution.
The output has the form $X_{\rm pre}+G/2$, where the reserved standard
Gaussian $G$ is independent of $X_{\rm pre}$ and all queries, and the
pre-reserved conclusion of Proposition~\ref{prop:compiler-accuracy}
gives $W_2(\mathcal L(X_{\rm pre}),\nu_{x,r}*N(0,3I_d/4))
\le\polylog(d)\delta^{K_0}D(x)$.
The exact means and numerical comparison laws serve only to analyze this finite
algorithm; its only oracle operations are the terminal evaluations.
\end{proof}
